\markdownRendererDocumentBegin
\markdownRendererSectionBegin
\markdownRendererHeadingOne{Discussion and conclusion}\markdownRendererInterblockSeparator
{}P12 reframes test-time scaling as a \markdownRendererStrongEmphasis{portfolio of computations}. “Think longer” is not the only adaptive action available to an intelligent system. It may be cheaper to parse, retrieve, compile, restructure or recover state so that less downstream search is required. Conversely, when state already exposes the relevant structure, additional preprocessing is wasteful and reasoning should receive the marginal budget.\markdownRendererInterblockSeparator
{}The protected benchmark establishes the key causal discriminator: when held-out families contain heterogeneous resource loci, a policy that can move the same total budget across both loci strictly dominates policies that may adapt only one. The result is not an artifact of giving the joint arm more computation, and it survives family-level uncertainty and worst-family reporting.\markdownRendererInterblockSeparator
{}This creates a concrete systems hypothesis for real agents: \markdownRendererStrongEmphasis{test-time scaling curves should be two-dimensional, with state-work and reasoning-work measured on a common receipt, and optimal allocation should shift with the source of difficulty.}\markdownRendererInterblockSeparator
{}Adaptive inference should decide not only \markdownRendererStrongEmphasis{how much} computation to spend but \markdownRendererStrongEmphasis{where} to spend it. P12 supplies a matched-budget formulation and a protected controlled result showing strict, broad held-out superiority of joint state–reasoning allocation over both one-axis adaptive alternatives. The next scientific step is not a larger synthetic benchmark; it is a real end-to-end system in which compiler/retrieval work and downstream reasoning/search are charged together and the same superiority contract is tested without privileged signals.
\markdownRendererSectionEnd \markdownRendererDocumentEnd